% Generated by tools/build_content_sets.py. Standard LaTeX; no extra packages.
\begin{table}[t]
\centering
\caption{Proposed target-tab positions in the high-load condition.}
\label{tab:target-positions}
\small
\begin{tabular}{lccccc}
\hline
Template (content set) & P1 & P2 & P3 & P4 & P5 \\
\hline
A (Sports) & 2 & 4 & 8 & 11 & 15 \\
B (Animals) & 1 & 6 & 9 & 12 & 14 \\
C (Nature) & 3 & 5 & 7 & 10 & 13 \\
D (Science) & 2 & 6 & 8 & 12 & 15 \\
\hline
\end{tabular}
\medskip
\begin{minipage}{\linewidth}
\footnotesize
\textit{Note.} P1--P5 identify target pages. Study Start is last: browser
tab 6 at low load or 16 at high load. Task position $j$ equals browser
tab $j$. Low load retains P1--P5 in positions 1--5; high load adds ten
distractors in predetermined positions. Tab order remains fixed within
each block. Each high-load template includes one unique target in each
band: 1--3, 4--6, 7--9, 10--12, and 13--15. Repeated-target trials add
another observation at the repeated target's position.
\end{minipage}
\end{table}

\begin{table}[t]
\centering
\caption{Proposed rotation of content sets and position templates across conditions.}
\label{tab:template-assignment}
\small
\begin{tabular}{ccccc}
\hline
Group & H--L & H--H & V--L & V--H \\
\hline
1 & A & B & C & D \\
2 & B & C & D & A \\
3 & C & D & A & B \\
4 & D & A & B & C \\
\hline
\end{tabular}
\medskip
\begin{minipage}{\linewidth}
\footnotesize
\textit{Note.} H--L = horizontal-low; H--H = horizontal-high;
V--L = vertical-low; V--H = vertical-high. Low load has five task
tabs plus Study Start; high load has fifteen plus Study Start. Allocate
participants equally across groups (four per group if $N=16$). This
table specifies material assignments, not chronological condition order.
Coordinate condition order and target-prompt schedules separately.
\end{minipage}
\end{table}
